
\documentclass{article}
\usepackage{colortbl}
\usepackage{makecell}
\usepackage{multirow}
\usepackage{supertabular}

\begin{document}

\newcounter{utterance}

\centering \large Interaction Transcript for game `cladder', experiment `full\_v1.5\_default', episode 7517 with reflexion{-}Qwen3.5{-}9B.
\vspace{24pt}

{ \footnotesize  \setcounter{utterance}{1}
\setlength{\tabcolsep}{0pt}
\begin{supertabular}{c@{$\;$}|p{.15\linewidth}@{}p{.15\linewidth}p{.15\linewidth}p{.15\linewidth}p{.15\linewidth}p{.15\linewidth}}
    \# & \multicolumn{2}{c}{Player} && \multicolumn{2}{c}{Game Master} \\
    \hline

    \theutterance \stepcounter{utterance}  
    & & & \multicolumn{4}{p{0.6\linewidth}}{
        \cellcolor[rgb]{0.9,0.9,0.9}{
            \makecell[{{p{\linewidth}}}]{
                \texttt{\tiny{[P1$\langle$GM]}}
                \texttt{You are an expert in causal inference. The following question is not a typical commonsense query, but rather a meticulously designed question created by a professor specializing in causal inference, intended to assess the students' mastery of the course content.} \\
\texttt{Imagine a self{-}contained, hypothetical world with only the following conditions, and without any unmentioned factors or causal relationships: Wibl has a direct effect on vubr and zuph. Zuph has a direct effect on uvzi. Vubr has a direct effect on uvzi. Method 1: We look at how zuph correlates with uvzi case by case according to vubr. Method 2: We look directly at how zuph correlates with uvzi in general. To understand how zuph affects uvzi, is it more correct to use the Method 1 than Method 2?} \\
\texttt{Start your answer with 'Yes' or 'No', followed by additional reasoning or evidence to support your explanation.} \\
            }
        }
    }
     \\ \\

    \theutterance \stepcounter{utterance}  
    & \multicolumn{4}{p{0.6\linewidth}}{
        \cellcolor[rgb]{0.561,0.737,0.561}{
            \makecell[{{p{\linewidth}}}]{
                \texttt{\tiny{[P1$\rangle$GM]}}
                \texttt{No} \\
\\ 
\texttt{**Reasoning:**} \\
\\ 
\texttt{In this hypothetical world, the causal structure forms a **confounding scenario** (specifically, a "fork" or common cause structure relative to the variables of interest, though here it is a collider{-}like dependency in terms of correlation patterns).} \\
\\ 
\texttt{1.  **Causal Graph Analysis**:} \\
\texttt{    *   \$Wibl \textbackslash{}rightarrow Zuph\$} \\
\texttt{    *   \$Wibl \textbackslash{}rightarrow Vubr\$} \\
\texttt{    *   \$Zuph \textbackslash{}rightarrow Uvzi\$} \\
\texttt{    *   \$Vubr \textbackslash{}rightarrow Uvzi\$} \\
\\ 
\texttt{    Here, \$Wibl\$ is a common cause of both \$Zuph\$ and \$Vubr\$. Both \$Zuph\$ and \$Vubr\$ are direct causes of \$Uvzi\$.} \\
\\ 
\texttt{2.  **Analyzing Method 2 (Unconditional Correlation)**:} \\
\texttt{    When you look directly at how \$Zuph\$ correlates with \$Uvzi\$ in general (Method 2), you are observing the marginal correlation. Because \$Wibl\$ influences both \$Zuph\$ and \$Vubr\$, and \$Vubr\$ influences \$Uvzi\$, the variation in \$Uvzi\$ is driven by a mix of \$Zuph\$ and \$Vubr\$. Since \$Zuph\$ and \$Vubr\$ are correlated (both caused by \$Wibl\$), the total correlation between \$Zuph\$ and \$Uvzi\$ will be a combination of the direct causal effect of \$Zuph\$ and the indirect association induced by their shared parent \$Wibl\$ (which also drives \$Vubr\$). This results in a biased estimate of the direct causal effect of \$Zuph\$ on \$Uvzi\$.} \\
\\ 
\texttt{3.  **Analyzing Method 1 (Stratification/Conditioning)**:} \\
\texttt{    Method 1 involves looking at the correlation between \$Zuph\$ and \$Uvzi\$ "case by case according to \$Vubr\$." This is equivalent to **stratifying** on \$Vubr\$ (or conditioning on it).} \\
\texttt{    *   By holding \$Vubr\$ constant (fixing the value of \$Vubr\$), you block the path through which \$Wibl\$ influences \$Uvzi\$ via \$Vubr\$.} \\
\texttt{    *   More importantly, in this specific DAG, \$Vubr\$ is a **confounder** of the relationship between \$Zuph\$ and \$Uvzi\$ because \$Wibl \textbackslash{}rightarrow Zuph \textbackslash{}rightarrow Uvzi\$ and \$Wibl \textbackslash{}rightarrow Vubr \textbackslash{}rightarrow Uvzi\$. The variable \$Vubr\$ is a common effect of \$Wibl\$ and a cause of \$Uvzi\$, but crucially, \$Zuph\$ and \$Vubr\$ are correlated due to \$Wibl\$.} \\
\texttt{    *   Actually, let's re{-}evaluate the} \\
            }
        }
    }
    & & \\ \\

    \theutterance \stepcounter{utterance}  
    & & & \multicolumn{2}{p{0.3\linewidth}}{
        \cellcolor[rgb]{0.9,0.9,0.9}{
            \makecell[{{p{\linewidth}}}]{
                \texttt{\tiny{[GM$|$GM]}}
                \texttt{No} \\
            }
        }
    }
    & & \\ \\

    \theutterance \stepcounter{utterance}  
    & & & \multicolumn{2}{p{0.3\linewidth}}{
        \cellcolor[rgb]{0.9,0.9,0.9}{
            \makecell[{{p{\linewidth}}}]{
                \texttt{\tiny{[GM$|$GM]}}
                \texttt{yes} \\
            }
        }
    }
    & & \\ \\

    \theutterance \stepcounter{utterance}  
    & & & \multicolumn{2}{p{0.3\linewidth}}{
        \cellcolor[rgb]{0.9,0.9,0.9}{
            \makecell[{{p{\linewidth}}}]{
                \texttt{\tiny{[GM$|$GM]}}
                \texttt{game\_result = LOSE} \\
            }
        }
    }
    & & \\ \\

    \theutterance \stepcounter{utterance}  
    & & \multicolumn{4}{p{0.6\linewidth}}{
        \cellcolor[rgb]{0.95,0.95,0.95}{
            \makecell[{{p{\linewidth}}}]{
                \texttt{\tiny{[P1$\leftrightarrow$P]}}
                \texttt{Post{-}episode Reflexion (not sent to GM):} \\
\texttt{New plan: I failed because I incorrectly identified the causal structure, mislabeling \$Vubr\$ as a confounder and \$Wibl\$ as a common cause of \$Zuph\$ and \$Vubr\$ when the prompt explicitly states \$Wibl\$ only affects \$vubr\$ and \$zuph\$ separately (no arrow between them). Consequently, I wrongly concluded that conditioning on \$Vubr\$ was necessary to block a backdoor path. In reality, since \$Wibl\$ is not a common cause of \$Zuph\$ and \$Vubr\$, there is no backdoor path between them; the only path is the direct causal chain \$Zuph \textbackslash{}rightarrow Uvzi\$. Therefore, Method 2 (unconditional correlation) correctly captures the direct effect, making Method 1 (conditioning on \$Vubr\$) incorrect as it would block the causal path. Next time, I will strictly map the arrows from the prompt first: \$Wibl \textbackslash{}to Zuph\$, \$Wibl \textbackslash{}to Vubr\$, \$Zuph \textbackslash{}to Uvzi\$, \$Vubr \textbackslash{}to Uvzi\$. I will then identify that \$Zuph\$ and \$Vubr\$ are independent given no common cause, meaning no adjustment is needed, and answer "No" with reasoning based on the absence of confounding.} \\
            }
        }
    }
    & \\ \\

\end{supertabular}
}

\end{document}
